\documentclass{beamer}
\def\presentationtype{0}
\def\corsodilaurea{Ingegneria Informatica e dell'Automazione Industriale}
\def\insegnamento{Ottimizzazione}
\def\nolezione{8}
\def\titolo{Il metodo del simplesso -- Il ``tableau'' e l'operazione ``pivot''}
\def\noattivita{1}

\def\argomenti{Programmazione Lineare}
\def\obiettivilezione{Definire il ``tableau'' e l'operazione ``pivot'' per poter sviluppare metodo del simplesso primale}
\def\nucleotematico{Richiami}
\def\contenuto{Programmazione lineare}


\usepackage[italian]{babel}
%
% Hyperref info for PDF
%
\usepackage{hyperref}
\hypersetup{%
%pdfpagemode=FullScreen,
pdfpagemode=FullScreen,
pdfstartview=Fit,
pdfsubject={\argomenti},
pdfkeywords={\insegnamento,\titolo}
}

\definecolor{ecampus@engred}{RGB}{227,28,23}
\setbeamercolor{normal text}{fg=black,bg=}
\setbeamercolor{alerted text}{fg=red}
\setbeamercolor{example text}{fg=green!50!black}
\setbeamercolor{structure}{fg=ecampus@engred,bg=}
\usetheme{default}
\useinnertheme{rounded}%{circles}
\setbeamertemplate{blocks}[default]
\setbeamercolor{block title}{bg=}
\setbeamercolor{block body}{bg=}

\ifnum\presentationtype>0 {\setbeamertemplate{enumerate item}{Es. \nolezione.\presentationtype.\arabic{enumi})}} \fi

%\setbeamertemplate{frametitle continuation}[from second][(cont'd)]
\usefonttheme[stillsansseriftext,stillsansserifsmall]{serif}
\setbeamertemplate{headline}{%
\leavevmode%
   \hspace*{4.75cm}
   \scalebox{0.90}{%
    \begin{minipage}{\linewidth}%
      \tiny
 	\begin{tabular}{ll}
	  \\
% 	  \textcolor{gray}{Classe}					& \corsodilaurea\\
%	  \textcolor{gray}{Argomento:}					& \insegnamento\\
%	  \textcolor{gray}{Lezione n}:							& \nolezione \\
	  \textcolor{gray}{Titolo:}			& \titolo \\
	  \textcolor{gray}{Attivit\`a}:		& 	  \ifnum\presentationtype=0 {Lezione \nolezione} \else {Sessione di studio \nolezione.\presentationtype} \fi
	\end{tabular}%
    \end{minipage}%
   }%
}
\setbeamertemplate{footline}[text line]{%
  \begin{minipage}{\linewidth}%
    \hfill\insertpagenumber%
%    
    \begin{center}%
      \rule{\linewidth}{0.4pt}
    \end{center}
    \vspace*{-0.65cm}
    \begin{center}%
      \scalebox{0.55}{%
	\begin{minipage}{\linewidth}%
	  \begin{center}%
	    {\tiny
	    \copyright \ 2022 Gionata Massi - \color{blue}{\href{mailto:gionata.massi@savoiabenincasa.it}{gionata.massi@savoiabenincasa.it}}\\
	    \phantom{XXX}}
	  \end{center}
	\end{minipage}%
      }%
    \end{center}%
%
  \end{minipage}%
  }
\setbeamertemplate{navigation symbols}{}

%% Save up on ink for the 4-up handouts
\mode<handout>{%
  \pgfpagesuselayout{4 on 1}[a4paper, landscape, border shrink=10mm]
  \pgfpageslogicalpageoptions{1}{border code=\pgfstroke}
  \pgfpageslogicalpageoptions{2}{border code=\pgfstroke}
  \pgfpageslogicalpageoptions{3}{border code=\pgfstroke}
  \pgfpageslogicalpageoptions{4}{border code=\pgfstroke}
}

\mode<presentation>{\AtBeginSection{%
  \begin{frame}
    \frametitle{Piano della presentazione}
    \tableofcontents[currentsection]
  \end{frame}}
}

\usepackage{microtype}
\usepackage[utf8]{inputenc}
\usepackage[T1]{fontenc}
%\usepackage[osfss]{libertine}
\usepackage{palatino}
\usepackage[scaled=.77]{beramono}
\usepackage{booktabs}
% \usepackage{attachfile}

% Toglie l'enorme spazione prima di description
%\setbeamersize{description width=0pt}

\usepackage{tikz}
\usetikzlibrary{decorations,arrows,shapes,backgrounds,matrix,positioning}
\usepackage{multirow}
\usepackage{verbatim}
\usepackage{eurosym}
\usepackage{pgfpages}

\usepackage{zref-savepos}
\newcounter{restofframe}
\newsavebox{\restofframebox}
\newlength{\mylowermargin}
\setlength{\mylowermargin}{2pt}

\newenvironment{restofframe}{%
    \par%\centering
    \stepcounter{restofframe}%
    \zsavepos{restofframe-\arabic{restofframe}-begin}%
    \begin{lrbox}{\restofframebox}%
}{%
    \end{lrbox}%
    \setkeys{Gin}{keepaspectratio}%
    \raisebox{\dimexpr-\height+\ht\strutbox\relax}[0pt][0pt]{%
    \resizebox*{!}{\dimexpr\zposy{restofframe-\arabic{restofframe}-begin}sp-\zposy{restofframe-\arabic{restofframe}-end}sp-\mylowermargin\relax}%
        {\usebox{\restofframebox}}%
    }%
    \vskip0pt plus 1filll\relax
    \mbox{\zsavepos{restofframe-\arabic{restofframe}-end}}%
    \par
}

\title{\titolo}
\author{Gionata Massi}
\date{\today}

\setbeamercolor*{frametitle}{fg=black}
\setbeamercolor*{title}{fg=black}
\setbeamerfont{frametitle}{shape=\scshape,family=\rmfamily,size=\large,series=\bfseries}

\setbeamersize{}
\makeatletter
\newcommand\thefontsize[1]{{#1 The current font size is: \f@size pt\par}}
\makeatother

\tikzstyle{nicebox}=[draw=gray!100, fill=blue!10, very thick,
rounded corners, rectangle, inner sep=4pt, inner ysep=16pt]
\tikzstyle{niceboxtitle}=[draw=gray!100, fill=white, text=black,
rounded corners, very thick, rectangle]
\newcommand\nicebox[2]{
{\centering
\begin{tikzpicture}
\node [nicebox](box){
\begin{minipage}{0.95\textwidth}\centering
\begin{minipage}{0.95\textwidth}
#2
\end{minipage}\end{minipage}};
\node[niceboxtitle, right=10pt] at (box.north west)
{\small\textbf{#1}};
\end{tikzpicture}\par}
}

\tikzstyle{modelbox}=[draw=structure!100, fill=white!50, very thick,
rounded corners, rectangle, inner sep=4pt, inner ysep=16pt, text=blue!50!black]
\tikzstyle{modelboxtitle}=[draw=structure!100, fill=white!50, text=blue!50!black,
rounded corners, very thick, rectangle]
\newcommand\modelbox[2]{
{\centering
\begin{tikzpicture}
\node [modelbox](box){
\begin{minipage}{0.95\textwidth}\centering
\begin{minipage}{0.95\textwidth}
#2
\end{minipage}\end{minipage}};
\node[modelboxtitle, right=10pt] at (box.north west)
{\small\textbf{\mbox{#1}}};
\end{tikzpicture}\par}
}

% % %
% Tipi di sessioni di studio
% % %
\def\domande{Domande}
\def\approfondimenti{Approfondimenti}
\def\esercizi{Esercizi}
\def\esempi{Esempi svolti}
\def\relazioni{Relazioni}
\def\attivitapratica{Attivit\`a pratica}
% % %

\newcommand{\generatitolo}{ %
\begin{frame}[plain]
\pdfbookmark{\ifnum\presentationtype=0 {Lezione \nolezione} \else {Sessione di studio \nolezione.\presentationtype} \fi}{\ifnum\presentationtype=0 {Lezione \nolezione} \else {Sessione di studio \nolezione.\presentationtype} \fi}
\noindent\hspace*{3.715cm}
   \scalebox{0.90}{%
    \begin{minipage}{\linewidth}%
      \tiny
 	\begin{tabular}{ll}
	  \\
% 	  \textcolor{gray}{Classe}					& \corsodilaurea\\
%	  \textcolor{gray}{Argomento:}					& \insegnamento\\
%	  \textcolor{gray}{Lezione n}:							& \nolezione \\
	  \textcolor{gray}{Titolo:}			& \titolo \\
	  \textcolor{gray}{Attivit\`a}:		& 	  \ifnum\presentationtype=0 {Lezione \nolezione} \else {Sessione di studio \nolezione.\presentationtype} \fi
	\end{tabular}%
    \end{minipage}%
   }%
   
  \begin{center}
    \large{\textbf{\insegnamento}}
  \end{center}

  \vspace*{0.5cm}
  
  \begin{center}
    \huge{\titolo}
  \end{center}
  
  \begin{center}
	  \ifnum\presentationtype=0
	    {Lezione n$^{\circ}$ \nolezione}
	  \else
	    {Sessione di studio n$^{\circ}$ \nolezione .\presentationtype}
	  \fi
  \end{center}
  
  \vspace*{0.5cm}

  \begin{center}
    \large{\textsl{Gionata Massi}}\\
    \scriptsize\color{blue}{<\href{mailto:gionata.massi@savoiabenincasa.it}{gionata.massi@savoiabenincasa.it}>}
  \end{center}
  
  \newcounter{pageleft}
  \zsaveposy{pageleft}
  \vspace*{\dimexpr\zposy{pageleft}sp-1.67cm}

  \begin{minipage}{\linewidth}%
   \hfill\phantom{\insertpagenumber}%

  \begin{center}%
      \rule{\linewidth}{0.4pt}
    \end{center}
    
    \vspace*{-0.95cm}
    
    \begin{center}%
      \scalebox{0.55}{%
	\begin{minipage}{\linewidth}%
	  %\begin{center}%
	  %  \tiny
	  %  \copyright \ 2022 Gionata Massi - \color{blue}{\href{mailto:gionata.massi@savoiabenincasa.it}{gionata.massi@savoiabenincasa.it}}\\
	  %  \color{blue}{\href{mailto:gionata.massi@savoiabenincasa.it}{gionata.massi@savoiabenincasa.it}}\\
	  %  \phantom{XXX}
	  %\end{center}
	\end{minipage}%
      }%
    \end{center}%
     \phantom{XXX}
%
  \end{minipage}%

\end{frame}
}

\renewcommand{\vec}{\mathbf}
\newcommand{\matr}{\mathbf}

\uselanguage{italian}
\languagepath{italian}
\deftranslation[to=italian]{Theorem}{Teorema}
\deftranslation[to=italian]{theorem}{teorema}
\deftranslation[to=italian]{Definition}{Definizione}
\deftranslation[to=italian]{definition}{definizione}
\deftranslation[to=italian]{Corollary}{Corollario}
\deftranslation[to=italian]{corollary}{corollario}

\let\definition\relax
\let\enddefinition\relax

%\theoremstyle{example}
\newtheorem{definition}[theorem]{\translate{Definition}}

\def\fieldN{\mathbb{N}}
\def\fieldR{\mathbb{R}}
\def\fieldZ{\mathbb{Z}}

\def\matrA{\matr{A}}
\def\vecB{\vec{b}}
\def\vecC{\vec{c}}
\def\vecX{\vec{x}}

\usepackage{cancel}
\usepackage{bookmark}

% tkiz ball item
\newcommand*\circled[1]{\raisebox{2.5pt}{\tikz[baseline=(char.base)]{
            \node[circle,ball color=purple, shade, 
 color=white,inner sep=1.2pt] (char) {\tiny #1};}}}

% tkiz rounded item
\newcommand*\rounded[1]{\tikz[baseline=(char.base)]{
            \node[draw=none,ball color=purple, shade, 
 color=white, rounded corners=3.5pt, inner sep=2.5pt] (char) {\scriptsize #1};}}

 \usepackage{etoolbox}
\makeatletter
\apptocmd{\beamer@@frametitle}{\write\@auxout{\string\@writefile{frm}{\string\frametitleentry{\the\c@framenumber}{#1}{#2}}}}{}{}
\newcommand*{\frametitleentry}[3]{\@namedef{frametitleshort#1}{#2}\@namedef{frametitle#1}{#3}}
\AtEndDocument{\if@filesw\newwrite\tf@frm\immediate\openout\tf@frm\jobname.frm\relax\fi}
\@input{\jobname.frm}
\makeatother
 
%\usebackgroundtemplate{\includegraphics[width=\paperwidth]{../img/sfondo_ecampus.png}}

\usepackage{listings}
% Extension for listings package
\lstdefinelanguage{ampl}{
  % See AMPL book (1993 edition), Appendix A
  morekeywords={%
    % Table A.1
    % Excluding symbols and "not X", where X (and not) are keywords
    if, then, else, or,  exists, forall,  and, in, within, 
    not, union, diff, symdiff, inter, cross, setof, by, 
    % Skipping: less,  sum, prod, min, max, div, mod,
    % Skipping built-in functions (Tables A-2 and A-3)
    % Section A.5 -- Declarations of model entities (subject to ->
    % subject, to)
    set, param, var, arc, minimize, maximize, subject, to, node,
    % Section A.6 -- Set declarations (not repeating set or within)
    dimen, default, card, ordered, by, reversed, circular,
    % Section A.7 -- Parameter declarations (not repeating in, default)
    param, binary, integer, symbolic, check, Infinity
    % Section A.8 -- Variable declarations
    var,  coeff, cover, obj,
    % Section A.13 -- Command language
    % Skipping most commands!
    include, model, data, let, objective, drop, restore, solve,
    solution, quit, end, option, reset, update, commands,
    % A.14 -- Imported functions
    function, pipe, 
    % A.17 synonyms for previously-defined keywords
    coef, difference, dimen, dimension, intersect, intersection, s.t., 
    % 
    % must be defined somewhere ...
    for, while, repeat, break,
    % NEW language features (http://www.ampl.com/NEW/newlang.html)
    % Grouped by page
    purge, problem, redeclare, xref,
    complements, 
    table, IN, OUT, INOUT, read, write, 
    expand,
    for, commands, repeat, until, while, break, continue,
  },
  sensitive=true,
  comment=[l]{\#},
  morecomment=[s]{/*}{*/},
  string=[b]",
  morestring=[b]',
}[keywords,comments,strings]

\definecolor{comment_c}{RGB}{60,128,49}
\definecolor{rowname_c}{RGB}{128,60,49}
\definecolor{reserved_c}{RGB}{49,60,128}

\lstset{
	language=ampl,%
	commentstyle=\color{comment_c},%
	backgroundcolor=\color{blue!15},%
    basicstyle=\small\ttfamily,%
    %emph={s.t.},emphstyle={\bfseries},%
    otherkeywords={s.t.},%
    numbers=left, numberstyle=\tiny, stepnumber=1, numbersep=5pt,%
}

\definecolor{col_set}{RGB}{128,0,0}
\definecolor{col_par}{RGB}{91,0,91}
\definecolor{col_idx}{RGB}{0,91,91}
\definecolor{col_var}{RGB}{0,0,128}

\def\insieme{%1}{\rm \textcolor{col_set}{%1}}}


\begin{document}

\generatitolo

\section{I fogli elettronici e la PL}

\begin{frame}{Fogli elettronici}
\begin{itemize}
\item I fogli elettronici (fra cui \emph{Microsoft Excel} e \emph{LibreOffice Calc}) sono
strumenti molto diffusi che permettono analizzare e risolvere
problemi di PL di piccole dimensioni.

\item Un foglio elettronico pu\`o essere utilizzato per:
\begin{itemize}
	\item visualizzare i dati;
	\item analizzare rapidamente le possibili soluzioni.
\end{itemize}

\item Il Risolutore incluso in \emph{Excel}, opportunamente configurato,
applica il metodo del simplesso per determinare una soluzione ottima
per l'instanza del modello di PL. Per altre tipologie di problemi di
ottimizzazione utilizza altri algoritmi.

\item \emph{LibreOffice} \emph{Calc} utilizza il metodo del
simplesso per risolvere problemi di PL.
\end{itemize}
\end{frame}

\begin{frame}{Programmi Lineari}
\begin{itemize}
\item Quando si inserisce un programma lineare su un foglio elettronico
\`e necessario conoscere i parametri del modello e distinguere:
\begin{enumerate}
\item le decisioni che devono essere prese (\alert{variabili decisionali});
\item le restrizioni sul valore che pu\`o essere assunto delle variabili
	decisionali (\alert{vincoli});
\item la misura di prestazione complessiva del sistema in funzione delle decisioni (\alert{funzione obiettivo})
\end{enumerate}
\end{itemize}
\end{frame}

\section{Esempio}

\begin{frame}[allowframebreaks]{Problema di esempio}
\fontsize{7}{8.4}{
\begin{block}{Mix ottimo di produzione}
Un'azienda di produzione vuole determinare il tasso di produzione mensile di due prodotti
in modo da massimizzare il profitto netto totale,  sapendo che:
 \begin{itemize}
  \item per produrre un quintale di prodotto 1 occorrono 40 quintali di materia prima e 8 ore di lavoro;
  \item per produrre un quintale di prodotto 2 occorrono 20 quintali di materia prima e 2 ore di lavoro;
  \item il commerciale ha stabilito che la produzione totale mensile non pu\`o superare 100 quintali;
  \item la disponibilit\`a mensile di materia prima \`e di 2200 quintali e quella di lavoro di 320 ore;
  \item il profitto netto per la vendita dei prodotti 1 e 2 sia rispettivamente 120 e 40 euro.
 \end{itemize}
\end{block}
}

\framebreak

\modelbox{Modello}{%
\small{%
$$\begin{array}{crrrlcr}
\max z=& 120 x_1 &+& 40 x_2    &           & (0) & [\text{profitto}]\\
{\rm s.t.} & 40 x_1 &+& 20 x_2 & \leq 2200 & (1) & [\text{materie prime}]\\
           &  8 x_1 &+&  2 x_2 & \leq 320  & (2) & [\text{lavoro}]\\
           &    x_1 &+&    x_2 & \leq 100  & (3) & [\text{mercato}]\\
           \multicolumn{5}{c}{x_1 \geq 0,\ x_2 \geq 0} & & [\text{non negativit\`a}]
\end{array}$$}}
\end{frame}

\section{Formulare il problema}

\begin{frame}{Foglio di calcolo}
\begin{itemize}
\item Nelle diapositive seguenti saranno presentate le schermate di
   \emph{LibreOffice Calc}, un foglio elettronico gratuito scaricabile da
   \href{http://it.libreoffice.org}{http://it.libreoffice.org},
   disponibile gratuitamente per numerose architetture hw/sw.
   
\item Le informazioni generali vanno bene per qualsiasi foglio elettronico
	mentre quelle sul risolutore sono specifiche per \emph{Calc}. % versione 4.
	
\item Per altri fogli elettronici controllare la documentazione appropriata.
\end{itemize}
\end{frame}


\begin{frame}{Inserire i parametri del problema}
\only<1>{
Per rappresentare un problema di PL servono
i suo elementi strutturali (variabili decisionali,
vincoli, funzione obiettivo) e i parametri del problema
specifico (istanza):

\begin{itemize}
\item un insieme di celle conterr\`a i valori delle variabili
	decisionali;
\item una cella del foglio elettronico conterr\`a il valore
	della funzione obiettivo calcolata in funzione del valore
	delle variabili decisionali;
\item un insieme di celle conterra i parametri del problema
	($\matrA,\ \vecB,\ \vecC,\ d$);
\item un gruppo di celle conterr\`a il valore calcolato dei
	valori dei termini di sinistra di ogni vincolo.
\end{itemize}
}
\includegraphics<2>[height=0.85\textheight]{./img/schermata1}
\end{frame}

\begin{frame}{Calcolare funzione obiettivo e primi membri}
\only<1>{
Abbiamo bisogno di calcolare la funzione obiettivo (lineare)
e il valore del primo membro di ogni vincolo lineare
(risultato dei prodotti scalari
$\vecC^T \vecX=z$ e $A_i \vecX \leq {\vecB}_i,\ i \in \{1, 2,\ldots,m\}$).

\begin{itemize}
\item \emph{Calc} ed \emph{Excel} offrono la funzione \emph{MATR.SOMMA.PRODOTTO}
	che permette di effettuare il prodotto scalare;
\item tale funzione vuole entrambi i vettori definiti come vettore riga
	o come vettore colonna;
\item per comodit\`a nel foglio di calcolo sono stati introdotti anche
	i valori delle variabili ausiliarie (slack).
\end{itemize}
}
\includegraphics<2>[height=0.85\textheight]{./img/schermata2}
\end{frame}

\begin{frame}{Verifica del corretto inserimento delle formule}
\only<1>{
\begin{itemize}
\item \`E consigliabile verificare di aver inserito le formule
correttamente nel foglio di calcolo;
\item un modo \`e quello di inserire dei valori
per le variabili decisionali e controllare che i risultati
siano quelli attesi.
\end{itemize}
}
\includegraphics<2>[height=0.85\textheight]{./img/schermata3}
\end{frame}

\section{Risolvere il problema}

\begin{frame}{Il risolutore}
\only<1>{
\begin{itemize}
\item Il risolutore \`e raggiungibile dal men\`u strumenti;
\item richiede di inserire i dati sul problema di PL:
\begin{itemize}
\item quale cella contiene il valore della funzione obiettivo;
\item la direzione di ottimizzazione ($\max$, $\min$ oppure
la minimizzazione dello scarto rispetto ad un certo valore);
\item le espressioni dei vincoli (primo termine, secondo termine,
	relazione tra loro).
\end{itemize}
\item La risoluzione di un modello di PL presuppone che tutte
	le variabili decisionali siano non negative.
\end{itemize}
}
\includegraphics<2>[height=0.85\textheight]{./img/schermata4}
\includegraphics<3>[height=0.85\textheight]{./img/schermata5}
\includegraphics<4>[height=0.85\textheight]{./img/schermata6}
\end{frame}


\begin{frame}{Il risolutore}
\only<1>{
\begin{itemize}
\item Il pulsante ``Risolvi'' permette di calcolare
	la soluzione ottima del problema utilizzando il
	metodo del simplesso.
\item Confrontando il valore della soluzione ottenuta
	col foglio di calcolo rispetto a quelle della
	risoluzione grafica e mediante il metodo del
	simplesso tableau notiamo che otteniamo lo stesso
	risultato:
	$\vecX^{T\star}=\left[\begin{array}{ccccc}25&60&0&0&15\end{array}\right]$
	in corrispondenza del quale $z^\star = 5400$.
	
\item Potremmo provare a modificare i valori per
	simulare decisioni diverse o variazioni dei
	parametri del modello.
\end{itemize}
}
\includegraphics<2>[height=0.85\textheight]{./img/schermata7}
\includegraphics<3>[height=0.85\textheight]{./img/schermata8}
\end{frame}

\section{Per approfondire}

\begin{frame}{Altro informazioni}
\begin{itemize}
\item Il risolutore di \emph{Excel} permette di effettuare
	anche l'analisi post-ottimale.
	
\item Esistono varie estensioni per aggiungere
	funzionalit\`a al risolutore di \emph{Calc}
	e di \emph{Excel}.
	
\item Entrambi i risolutori possono lavorare anche con
variabili intere (PLI e PLB).
	
\item \emph{Excel} e LibreoOffice con l'estensione
	\emph{nlpsolver} permettono di risolvere
	problemi non lineari.
\end{itemize}
\end{frame}

\begin{frame}{Per approfondire}
\beamertemplatebookbibitems
\begin{thebibliography}{References}
\bibitem{HillierLiebermanS1C3}F.S. Hillier, G. J. Lieberman.\newblock
\textit{Ricerca Operativa}. Capitolo 2.\newblock McGraw Hill, 2010.

\bibitem{ooo3solverug}\href{https://wiki.openoffice.org/wiki/Documentation/OOo3/_User/_Guides/Calc/_Guide/Solver}{Using the solver} \newblock \textit{OpenOffice.org} \newblock Pagina wiki consultata il 14 aprile 2022.

%\bibitem{Schrage2006} L. Schrage.\newblock \emph{Optimization Modeling
%with LINGO.} \newblock LINDO Systems, Inc., 2006

%\bibitem{LingoUserManual} LINDO Systems Inc.\newblock \emph{LINGO

%User's guide.} \newblock LINDO Systems, Inc., 2011
\end{thebibliography}
\end{frame}

\end{document}
